%% kapa: exam=1 round=30 year=2019 period=1 form=A subject=civil_law subject_ko=민법 seq=29 number=29
%% source: Q-Net 감정평가사 기출문제 (cst00309) artlSeq=5205205 file=extracted/1008/5205205/2019년 제30회 감정평가사 1차시험 문제지/2019년 제30회 감정평가사 1차시험 문제지/2019년도 제30회 감정평가사 1차 1교시 A형.hwp
%% conversion: kordoc 4.24.51 -> tools/gen1.py (auto); review=proofread-sonnet; flags=none
\begin{question}{29}
점유권의 효력에 관한 설명으로 옳지 않은 것은? (다툼이 있으면 판례에 따름)
\begin{choices}
  \item 점유자가 점유물에 대하여 행사하는 권리는 적법하게 보유한 것으로 추정된다.
  \item 점유자가 점유의 침탈을 당한 때에는 그 물건의 반환 및 손해의 배상을 청구할 수 있다.
  \item 점유물반환청구권은 점유의 침탈을 당한 날로부터 3년 내에 행사하여야 한다.
  \item 점유가 점유침탈 이외의 방법으로 침해되고 있는 경우에 점유자는 그 방해의 제거 및 손해의 배상을 청구할 수 있다.
  \item 점유권에 기인한 소와 본권에 기인한 소는 서로 영향을 미치지 아니한다.
\end{choices}
\end{question}
\answer{3}
